\section{Inne wybrane podejścia}

Poza ewolucją i metodami przeszukiwania drzewa gry, uczestnicy konkursu stosowali szereg innych technik dla swoich agentów.

\subsection{Alpha-Beta i Minimax}

Klasyczne podejście z teorii gier --- przeszukiwanie Minimax z przycinaniem Alfa-Beta --- zostało zaadaptowane do Hearthstone przez Wulfmeyera, który zajął 5. miejsce w edycji 2018 \cite{hearthstone_competition}. Wyzwaniem jest tu losowość mechanik gry: Minimax zakłada deterministyczność, więc losowe efekty kart wymagają uśredniania (ang.~\textit{expectimax}) lub przybliżenia deterministycznego. W edycji 2019 Sparks i Djartov posłużyli się agentem opartym na Min-Max, a w 2020 bot Lamprechta Mini-Max Search zajął 12. miejsce~\cite{hearthstone_competition}.

\subsection{Beam Search i BFS}

Beam Search --- uproszczenie przeszukiwania wszerz zachowujące tylko $k$ najobiecujących ścieżek --- pojawił się w edycji 2019 w bocie Bokelmanna i Unkringa oraz w edycji 2020 w bocie Ortlama. Pokrewnym podejściem jest BFS stosowane przez Heberta i Hecka w edycji 2019. Podejścia te redukują złożoność pamięciową przeszukiwania kosztem utraty gwarancji optymalności~\cite{hearthstone_competition}.

\subsection{Sieci neuronowe i uczenie maszynowe}

Kilku uczestników stosowało podejścia oparte na uczeniu maszynowym. Grad~\cite{Grad2017} zaproponował użycie sieci neuronowej jako zamiennika ręcznie zaprojektowanej heurystyki --- sieć była trenowana na danych z rozgrywek i służyła do oceny wartości stanu gry. Świechowski i in.~\cite{swiechowski2018} zastosowali modele nadzorowane jako sieć przewodnika w MCTS, kierując ekspansją drzewa ku akcjom obserwowanym u ludzkich graczy. W edycji 2019 Geis, Kleespies i Saflik zastosowali NN-based Tree Search w ścieżce \textit{User Created Deck Playing Track}~\cite{hearthstone_competition}.

\subsection{Symboliczne wnioskowanie}

Stiegler i in.~\cite{Stiegler2018} zaproponowali agenta opartego na wnioskowaniu symbolicznym, który koduje reguły strategiczne Hearthstone w postaci jawnych zasad logicznych. Podejście to kontrastuje z metodami czarnoskrzynkowymi: agent jest w pełni interpretowalny, lecz wymaga ręcznego specyfikowania wiedzy o grze i słabo skaluje się wraz ze wzrostem jej złożoności.

\subsection{Rolling Horizon Evolutionary Algorithm}

W edycji 2018 Weimer (15. miejsce, kat. \textit{Premade Deck Playing Track}) zastosował podejście Rolling Horizon Evolutionary Algorithm (RHEA). W tej metodzie, opisanej przez Gainę i in.~\cite{gaina2020}, algorytm ewolucyjny optymalizuje sekwencję akcji na bieżącą turę w czasie podejmowania decyzji --- bez wstępnego trenowania. W każdej chwili decyzji uruchamiana jest krótka sesja ewolucyjna, która wyznacza najlepszą sekwencję ruchów w ramach dostępnego czasu. Podejście to stanowi interesującą hybrydę między algorytmami \textit{online}, jak MCTS, a ewolucyjnymi, jednak w edycji 2018 nie dało wyników porównywalnych z czołowymi botami~\cite{hearthstone_competition}.